\section{Discussion}
\label{sec:discussion}

\subsection{Simplicity Outperforms Sophistication}
\label{sec:discussion-simplicity}

\texttt{tabpfn\_threshold} does less than any other correction tested. It
shifts the decision boundary from 0.5 to the class prior and nothing more
(Section~\ref{sec:methods-corrections}). Despite that, it is the most
robust method in the study on discrimination. It has the lowest CHR and
the highest mean CG in both directions, and its calibration is identical
to its own baseline's by construction, since a threshold shift never
touches the underlying predicted probabilities
(Section~\ref{sec:discussion-calibration} returns to why that is not the
same as leading on calibration outright). The one place it does not lead
on discrimination is break-even share, where \texttt{tabpfn\_downsample}
and \texttt{xgb\_ros} tie or exceed it under Direction A, and
\texttt{xgb\_smote} stays net-beneficial longer under Direction B, a
reminder that no single measure ranks every method the same way.
\texttt{tabpfn\_distpfn} and \texttt{tabpfn\_distpfn\_t} sit at the
opposite end on discrimination. Both rescale TabPFN's predictions against
a prior and a temperature term, and \texttt{tabpfn\_distpfn} has the
lowest mean CG and the highest regret of any method in both directions,
with \texttt{tabpfn\_distpfn\_t} close behind on both.

This finding builds on prior work rather than standing alone. Elor and
Averbuch-Elor~\cite{elor2022smote} found that on clean, uncorrupted data, a
well-tuned classifier often gains nothing from balancing the training set,
and can lose performance from it. This study's contribution is to show
that the same pattern holds, and grows sharper, once label noise is added.
A correction is not just sometimes unnecessary. It can actively backfire,
and in this comparison the two most mathematically elaborate corrections,
DistPFN and DistPFN-T, backfire hardest on discrimination, while the
simplest correction backfires least. A practitioner choosing a correction
method for a TabPFN deployment under uncertain label quality and
optimizing for balanced accuracy has, based on this study's results, no
clear reason to prefer a more elaborate posterior-rescaling method over a
simple threshold shift. That recommendation does not extend to every
deployment, however, a point Section~\ref{sec:discussion-calibration}
returns to.

\subsection{Direction Asymmetry Undercuts a Single Robustness Score}
\label{sec:discussion-direction}

Five of seven methods fail more often under Direction A (majority rows
relabeled as minority) than Direction B (minority rows relabeled as
majority), while the two DistPFN variants show the opposite pattern. This
asymmetry is exactly what the label-noise literature's class-conditional
and asymmetric-noise formalizations
predict in principle~\cite{natarajan2013noisylabels,scott2013asymmetricnoise},
but this study's results show it is not a minor theoretical caveat. A
method's CHR under one direction is a poor predictor of its CHR under the
other. \texttt{tabpfn\_downsample}'s CHR is 21.7\% under Direction A and
7.5\% under Direction B, nearly a threefold difference in the same method
on the same datasets. A robustness claim validated under only one
corruption mechanism, or under the symmetric random-flip noise more
commonly studied, does not transfer to a directed, class-conditional error
of the kind this study injects, and in practice the true direction of
label error is rarely known in advance. This is the strongest argument in
this study for testing every correction under both directions rather than
one.

\subsection{The Corrupted Prior Is Not the Main Failure Mode}
\label{sec:discussion-prior}

The most surprising result of this study is that mean PSG is negative for
all three prior-dependent TabPFN methods in both directions. The true,
uncorrupted prior, something no real deployment could ever supply, does
not reliably outperform the corrupted, observed one. Under Direction A,
at high contamination, the oracle variant of \texttt{tabpfn\_distpfn} and
\texttt{tabpfn\_distpfn\_t} actually degrades faster than the deployable
one. If the wrong prior were the primary cause of these methods'
decline, correcting it should have helped. It does not.

One structural reason follows directly from how these corrections work. \texttt{tabpfn\_threshold}, \texttt{tabpfn\_distpfn}, and
\texttt{tabpfn\_distpfn\_t} all adjust TabPFN's output probabilities after
the fact, using the class prior~\cite{saerens2002adjusting,lee2026distpfn}.
None of them touch the training rows TabPFN actually reads as its
in-context prompt~\cite{hollmann2023tabpfn,hollmann2025tabpfnv2}.
Contamination happens before any correction is applied, so the mislabeled
rows stay in that prompt regardless of which prior a posterior-rescaling
method is given afterward. Swapping in the true prior fixes the summary
statistic the correction uses, but not the individual wrong examples
TabPFN has already been shown. This would explain why the oracle prior
helps so little, though confirming it directly would require an
experiment this study did not run, comparing a prior-only fix against a
fix that also removes the mislabeled rows from the context itself.

\subsection{Accuracy and Calibration Answer Different Questions}
\label{sec:discussion-calibration}

Every XGBoost correction improves calibration over its baseline at every
noise level tested, even the ones that sometimes hurt balanced accuracy.
Under Direction A, \texttt{tabpfn\_downsample} is worse-calibrated than
doing nothing at every noise level tested. \texttt{tabpfn\_distpfn} and
\texttt{tabpfn\_distpfn\_t} take the opposite path entirely. The same two
methods with the highest Correction Harm Rate in the whole study are also
its two best-calibrated methods, correction or baseline, at nearly every
noise level in both directions. A method that frequently classifies a row
wrong can still be assigning that wrong classification a well-calibrated
probability. Balanced accuracy and calibration measure different
properties of a classifier, and this study's results show a correction
can excel at one while failing badly at the other, not merely land
somewhere middling on both. This matters for deployment. A system that
only needs a correct label, such as most classification benchmarks,
cares primarily about balanced accuracy or a related discrimination
metric, and by that standard DistPFN and DistPFN-T are a poor choice
under contamination. A system that acts on a predicted probability
(setting a risk threshold, allocating a limited budget, ranking cases by
urgency) cares about calibration first, and for that use case their poor
CHR would be the wrong reason to rule them out. Reporting CHR and mean CG
alone for the headline comparisons would also miss that
\texttt{tabpfn\_downsample} is uniquely bad at calibration under
Direction A even though its balanced accuracy profile looks far more
favorable there, the second-highest mean CG of any method. Evaluating a
correction method for deployment requires checking both axes separately,
not inferring one from the other, and not assuming a method's rank on one
predicts its rank on the other.

\subsection{Implications for Imbalance Correction Under Label Noise}
\label{sec:discussion-implications}

Related work on class imbalance and label noise as a joint problem has so
far proposed new training-time algorithms, combining cleaning with
resampling~\cite{koziarski2020ccr}, active learning with noisy-label
training~\cite{khanal2024activelabel}, or class-balanced sample
selection~\cite{liu2024imbalancednoisy}. This study starts from a
different question. It asks whether the corrections practitioners are
already using need replacing at all. The answer depends on which
correction is already in use. \texttt{tabpfn\_threshold} stays
net-beneficial through the full contamination range in half or more of
the conditions tested, in both directions, so switching to it may fix
most of what a new algorithm would otherwise need to solve.
\texttt{tabpfn\_distpfn} and \texttt{tabpfn\_distpfn\_t} fail badly
enough on discrimination under high contamination that the case for
joint noise-and-imbalance-aware training remains strong, for a
deployment that needs a correct label. A deployment that instead needs a
well-calibrated probability is a different case entirely
(Section~\ref{sec:discussion-calibration}), and swapping either DistPFN
variant out for a different correction there could trade away its best
property to fix its worst one. Telling these situations apart requires
the kind of audit this study performs, not an assumption that a
published correction is safe by default, or unsafe by default, on every
axis a deployment might care about.
